\section*{Hoofdstuk \ref{ch:g6:describing-motion} --- Beweging beschrijven: baan en snelheid}

\begin{solution}{exo:g6:describing-motion:1}
Rechtlijnig (langs de ruit omlaag); cirkelvormig; krom (een boog);
krom (een slingerende S-lijn).
\end{solution}

\begin{solution}{exo:g6:describing-motion:2}
De naaf: een rechte lijn, evenwijdig aan de weg glijdend. Het ventiel:
een kromme van sierlijke sprongetjes --- opstijgend, bogend, bij elke
wielomwenteling weer neerkomend.
\end{solution}

\begin{solution}{exo:g6:describing-motion:3}
Eenparig; veranderlijk (sneller wordend); eenparig; veranderlijk
(langzamer wordend).
\end{solution}

\begin{solution}{exo:g6:describing-motion:4}
Gelijke gaten: \omterm{def:g6:describing-motion:uniform}{eenparige beweging}. Groeiende gaten ($1, 2, 4, 7$):
sneller wordend --- meer afstand in elke gelijke tik.
\end{solution}

\begin{solution}{exo:g6:describing-motion:5}
Een afdwalende camera voegt zijn eigen beweging aan de opname toe ---
de stippen mengen dan de reis van het voorwerp met die van de camera.
Beweging is altijd beweging ten opzichte van iets, en de film moet dat
``iets'' vastleggen: de opmerking over gezichtspunten.
\end{solution}

\begin{solution}{exo:g6:describing-motion:6}
Tabel: $30$ min $\to$ \qty{12}{km}, dus $60$ min $\to$ \qty{24}{km}:
de \omterm{def:g5:speed-distance-time:speed}{snelheid} van de veerboot is $24$ kilometer per uur.
\end{solution}

\begin{solution}{exo:g6:describing-motion:7}
$60$ min $\to$ \qty{14}{km}; $30$ min $\to$ \qty{7}{km}; $15$ min
$\to$ \qty{3.5}{km}.
\end{solution}

\begin{solution}{exo:g6:describing-motion:8}
Een slenteraar met $12$ kilometer per uur: afgewezen --- drie keer het
slentertempo. De bus: $10$ km in $30$ min is $20$ kilometer per uur
--- precies stadsbustempo, aanvaard. De trein: $75$ km in $15$ min is
$300$ kilometer per uur --- bovenaan het bord maar geloofwaardig voor
een hogesnelheidstrein: aanvaard.
\end{solution}

\begin{solution}{exo:g6:describing-motion:9}
Voor de passagier tegenover hem: helemaal geen beweging --- geen
\omterm{def:g6:describing-motion:trajectory}{baan}, tempo nul. Voor de koe: een lange
\omterm{def:g6:describing-motion:trajectory}{rechtlijnige baan} met het volle tempo van de trein. Allebei
eerlijk: elk beschrijft de dommelaar ten opzichte van het eigen
gezichtspunt, en beweging is altijd beweging ten opzichte van iets.
\end{solution}

\begin{solution}{exo:g6:describing-motion:10}
Rijtijd: $18 - 6 = 12$ minuten voor \qty{9}{km}. Tabel: $12$ min
$\to$ $9$ km, dus $60$ min $\to$ $45$ km --- precies $45$ kilometer
per uur: niet erboven en niet eronder, maar er precies op.
\end{solution}

\begin{solution}{exo:g6:describing-motion:11}
Film het apporteren met een stilstaande camera; pauzeer elke halve
\omterm{def:g2:measuring-time:second}{seconde}; markeer de plaatsen van de hond ten opzichte van hekpalen;
meet de gaten. Overal gelijke gaten: gelijkmatig tempo. Gaten die
groeien en dan krimpen (of afwisselend breed en smal): sprinten en
uitrollen. De stippen vellen het oordeel, wat de familie ook riep.
\end{solution}

\begin{solution}{exo:g6:describing-motion:12}
Vanuit de kamer: een kromme --- de bal boogt van de tafelrand vooruit
en omlaag naar de vloer. Vanaf het karretje dat met het voorwaartse
tempo van de bal meerolt: de bal loopt niet voor en niet achter ---
hij lijkt \emph{recht omlaag} te vallen. Dezelfde val, twee
gezichtspunten, en het eenvoudigste \omterm{prop:g5:mirrors-reflection:image}{beeld} hoort bij de bewegende
waarnemer: een zaadje van een groot idee.
\end{solution}

\begin{solution}{pb:g6:describing-motion:1}
\textbf{1.} Gelijke gaten van \qty{50}{m}: \omterm{def:g6:describing-motion:uniform}{eenparige beweging}.
\textbf{2.} $5 \times 50 = \qty{250}{m}$.
\textbf{3.} De tram trekt op uit stilstand: korte gaten die groeien
--- sneller wordend over de eerste drie tikken --- en dan tot vaste
gaten van \qty{50}{m} bedaren: kruisen.
\textbf{4.} De laatste twee intervallen, elk \qty{50}{m}.
\textbf{5.} $10$ s $\to$ \qty{50}{m}, dus $60$ s $\to$ \qty{300}{m}
per minuut.
\textbf{6.} $60$ min $\to$ $60 \times 300 = \qty{18000}{m}$:
kruissnelheid $18$ kilometer per uur.
\textbf{7.} Tabel: \qty{18}{km} per $60$ min, dus \qty{12}{km} in
$40$ minuten.
\textbf{8.} $50 - 40 = 10$ minuten verloren aan haltes en
langzaamrijzones.
\textbf{9.} $50$ min $\to$ \qty{12}{km}, dus $60$ min $\to$
$12 \times 60 \div 50 = \qty{14.4}{km}$: de hele reis loopt op
ongeveer $14$ kilometer per uur --- de $18$ van de poster is alleen
eerlijk voor de kruisstukken, niet voor ``dwars door de stad''.
\textbf{10.} Een reissnelheid beschrijft de rit \emph{alsof} die
eenparig was: het ene gelijkmatige tempo dat dezelfde \qty{12}{km} in
dezelfde $50$ minuten zou afleggen. Geen snelheidsmeter toont haar, en
toch is ze het eerlijke getal om mee te plannen, omdat ze elke halte
en elke langzaamrijzone al bevat.
\textbf{11.} $50 - 4 = 46$ minuten;
$12 \times 60 \div 46 \approx \qty{15.7}{km}$ per uur --- zeg
ongeveer $16$.
\textbf{12.} Bijvoorbeeld: ``Reizigers moeten plannen met de
reissnelheid --- vandaag ongeveer $14$ kilometer per uur --- want dat
is het enige getal dat het wachten meetelt; de kruissnelheid van $18$
vleit de lijn en slaat elke halte over.''
\end{solution}
